% SPDX-License-Identifier: Apache-2.0
% covers: Board
\datasheet{Board}{The Vreteno machine for the Alinx AX7A200B, as one lowered unit of units.}

\dsfacts{\code{cpu/vreteno/src/board.rs}}{\code{vreteno32::board::Board}}%
{\code{//docs:vreteno}}

\subsection*{Function}

\code{Board} is everything that computes on the board, written as a
unit of units so that \code{\#[lower]} makes its netlist as one module
holding the rest. It holds seven hosts, an arbiter that puts them onto
one link, a router of eight, and the peripherals behind its eight
ports. The hosts are the core's tracker; Vivado's JTAG-to-AXI master
on the board's top (issue 241), whose pins arrive under \code{jtag\_}
and are joined by an \code{AxiPins}, sharing one host through
\code{jarb} with the bridge of the RISC-V debug transport (issue 154);
the Ethernet port's two engines, each behind an \code{AxiHost} of its
own (issue 151); the video scanout's fetch (issue 151); and the SD card
host's two engines, sharing one host through \code{sdarb} (issue 912);
and Razboj's rasteriser behind \code{rhost} (issue 985).
The arbiter is an \code{Arbiter7<32, 32, 4, 2, 5, 0>}, so the
peripheral side of the bus allocates five bits of identifier, two for a
host's own and three for the port, which is room for eight hosts
(issue 1022). The debug transport, \code{dtm}, runs on the cable's
clock behind the top's \code{BSCANE2}, and \code{dreq} and
\code{dans} transfer its accesses across to the board's clock and back.
The data memory and the timer each sit behind an \code{AxiPer}. Ten
small peripherals share the page at \code{0x3000} behind the AXI-Lite
bridge \code{LiteBridge<10, ..>}, a sixteenth of the page each: the
serial port at
\code{0x3000}, the pulse width modulator at \code{0x3100}, whatever a
board hangs on the slot at \code{0x3200}, the remote peripheral at
\code{0x3300}, the Ethernet port's registers at \code{0x3400}, the
entropy source at \code{0x3500} (issue 458), the configuration
flash's SPI master at \code{0x3600} (issue 312), the Ethernet PHY's
management interface at \code{0x3700} (issue 864), the SD card host
at \code{0x3800} (issue 153), and Razboj's doorbell at \code{0x3900}
(issue 985). The DDR3 memory takes the router's port
with no tracker in front, and holds AMD's MIG controller with its own
AXI4 port as a foreign module, which makes the design's clock from the
board's. The interrupt
controller, \code{Plic3}, and the debug module, \code{Dm}, each sit
behind a \code{LiteBridge<1, ..>} of their own; the module's lines go to
the core's halt and resume requests and its register access, and take
the core's debug mode and the word read (issue 154).

The configuration flash, the chip that holds the bitstream, is reached
twice (issue 312): through that SPI master, \code{spi}, for commands,
and through \code{flashwin}, a \code{FlashWin<FLASH\_DIV, 5>} behind
its own \code{AxiPer}, where an ordinary read of the 16~MiB from
\code{0x2000\_0000} is answered with what the flash holds.
\code{cfgflash}, a \code{CfgFlash<1>}, stands between both and the
chip: it reaches the clock pin through \code{STARTUPE2}, waits for the
end of configuration, holds the window in reset until then, and
refuses write enable from the master.

\code{run.rs} wires the same parts for a simulation. \code{Board} is
that wiring as a unit, so a board top has to supply only the clocks,
the resets and the pins. The top for the AX7A200B is
\code{cpu/vreteno/board/vreteno\_board.v}.

\subsection*{Parameters}

\begin{center}\footnotesize
\begin{tabular}{@{}lp{0.7\textwidth}@{}}
\toprule
Parameter & Meaning \\
\midrule
\code{DIV} & The serial port's cycles per bit, as \code{Uart<DIV>}.
The tables use 868, the value \code{board\_netlist board} writes for
115200 baud at 100~MHz. \\
\bottomrule
\end{tabular}
\end{center}

\code{board\_netlist} lowers \code{Board<868>} for the board and
\code{Board<16>} for simulation. \code{//cpu/vreteno:board\_test}
runs \code{Board<4>}.

\subsection*{Address map}

\code{BoardRouter} is \code{Router<8, BoardMap, 32, 32, 4, 5>}, and
\code{BoardMap} beside it in \code{board.rs} states the map. A
peripheral gets an address when the address under the mask equals the
base, the first range that matches deciding.

\begin{center}\footnotesize
\begin{tabular}{@{}lllp{0.35\textwidth}@{}}
\toprule
Base & Mask & Child & Behind \\
\midrule
\code{0x0000\_0000} & \code{0xffff\_f000} & \code{rom}, \code{Rom<5>} &
\code{prom}, \code{AxiPer}: the boot memory, readable and not writable \\
\code{0x1000} & \code{0xffff\_f000} & \code{dmem}, \code{Dmem<5>} &
\code{pdmem}, \code{AxiPer} \\
\code{0x0200\_0000} & \code{0xffff\_0000} & \code{timer},
\code{Timer<5>} & \code{ptimer}, \code{AxiPer} \\
\code{0x3000} & \code{0xffff\_f000} & \code{uart}, \code{Uart<DIV>} &
\code{puart}, \code{LiteBridge<10, ..>} at \code{0x3000} \\
& \code{0xffff\_ff00} & \code{pwm}, \code{Pwm} &
the same bridge, at \code{0x3100} \\
& \code{0xffff\_ff00} & the board's, on its own clock &
the same bridge, at \code{0x3200} \\
& \code{0xffff\_ff00} & \code{remote},
\code{Remote<REMOTE\_WAIT>} & the same bridge, at \code{0x3300} \\
& \code{0xffff\_ff00} & \code{eth}, \code{EthSlots} &
the same bridge, at \code{0x3400} \\
& \code{0xffff\_ff00} & \code{entropy}, \code{Entropy} &
the same bridge, at \code{0x3500} \\
& \code{0xffff\_ff00} & \code{spi}, \code{Spi} &
the same bridge, at \code{0x3600} \\
& \code{0xffff\_ff00} & \code{mdio}, \code{Mdio} &
the same bridge, at \code{0x3700} \\
& \code{0xffff\_ff00} & \code{sd}, \code{Sd} &
the same bridge, at \code{0x3800} \\
& \code{0xffff\_ff00} & \code{doorbell}, \code{Doorbell} &
the same bridge, at \code{0x3900} \\
\code{0x4000\_0000} & \code{0xc000\_0000} & \code{ddr3},
\code{Ddr3Per} & nothing: the router's port goes to it \\
\code{0x0c00\_0000} & \code{0xfc00\_0000} & \code{plic},
\code{Plic3<0>} & \code{pplic}, \code{LiteBridge<1, ..>} at
\code{0x0c00\_0000} \\
\code{0x1000\_0000} & \code{0xffff\_0000} & \code{dmod}, \code{Dm} &
\code{pdm}, \code{LiteBridge<1, ..>} at \code{0x1000\_0000}; not nearer the
interrupt controller, whose 64~MiB window the router would match first \\
\code{0x2000\_0000} & \code{0xff00\_0000} & \code{flashwin},
\code{FlashWin<FLASH\_DIV, 5>} & \code{pflash}, \code{AxiPer}: the
configuration flash, read only \\
\bottomrule
\end{tabular}
\end{center}

The router answers a burst that matches no range with \code{DecErr}.
The boot memory is on the bus at zero as well (issue 268), readable and
not writable; the core fetches from its own copy below \code{0x1000}
and from the bus above it.

\dstables{Board}

\subsection*{Behaviour}

\begin{itemize}
\item The core's tracker is \code{AxiHost<32, 32, 4, 2, 4>}: 32-bit
addresses and words, four lanes, two-bit identifiers, four of them.
\item \code{rst} goes to the core, the timer, the serial port and the
interrupt controller. The DDR3 controller has its own reset,
\code{sys\_rst}, active high, and hands back \code{ui\_rst}, high
until the clock it makes is good.
\item The interrupt controller's source 1 is the serial port's
interrupt line, its source 2 is \code{irq}, and its source 3 is the
Ethernet port's arrival; all three ask while high.
Its line is the core's external interrupt. The timer's line goes to
the core's \code{tirq}.
\item Of the core's outputs, only \code{halt} is a port. \code{instr}
and the \code{Writeback} go nowhere; \code{dbg} and \code{dbg\_rdata}
go to the debug module, whose requests and register access come back
to the core.
\item The join runs the timer before the core, since the core reads
the timer's line in the same step, and the serial port before the
interrupt controller before the core, for the same reason. \code{Ddr3Per} runs its controller
before its pins part for the same reason.
\item \code{calib} and the memory's pins are \code{Ddr3Per}'s,
brought out as the unit's own ports.
\item The slot at \code{0x3200} does not reach a field. It leaves as
five channel ports, \code{vaw}, \code{var} and \code{vw} out and
\code{vb} and \code{vr} back, because what sits there runs on a clock
of its own and the crossing is the board top's business. A board with
nothing there ties them off, and a program that touched \code{0x3200}
on such a board would wait for ever.
\item The remote peripheral at \code{0x3300} is the other way about:
\code{remote} and \code{link} are fields, because both run on the
core's clock. What leaves the unit is the frame stream, \code{net\_tx}
out and \code{net\_rx} in, a byte and a last bit each, which is what
the MAC speaks, and which the link shares with the Ethernet port
through \code{share}. \code{REMOTE\_DEV} is 1 and \code{REMOTE\_WAIT} is
100 million cycles, one second at 100~MHz (issue 397), after which an
unanswered transaction is \code{SLVERR} on the bus.
\item \textbf{The Ethernet port moves frames itself.} \code{eth} is
the register map; behind it the fetch engine reads a frame out of its
transmit slot and \code{FrameOut} turns the words into bytes, and on
the way in \code{FrameLen} finds the frame's length again,
\code{FrameIn} packs its bytes into words and the store engine writes
them into a receive slot. The slots are in the DDR3 memory from
\code{0x4100\_0000}. Frames the driver has not acknowledged wait in
order, at most one a slot, and while both slots hold one \code{eth}
tells \code{fin}, which drops the next frame and counts it for
\code{rx\_errors} (issue 1313). \code{NoBeats} and \code{NoReads} hold the
direction of each engine's host that the engine never uses.
\item \textbf{The wire is shared by EtherType.} \code{share}, an
\code{EthShare<0x88b5>}, sends the remote peripheral's frames to its
link and every other frame to the Ethernet port, and merges the two
users' frames onto \code{net\_tx} a frame at a time.
\item The join runs the receive side, \code{share}, \code{flen} and
\code{fin}, before the register block, because the block reads the
slot and length they drive in the same step; and the store engine and
the sending side after it, because they read what it drives. With the
receive side after the block, the block read last step's slot, and
the board test caught it as a frame of the right length whose every
byte read back zero. \code{fin} reads the block's \code{rx\_full} a
step late in the Rust run for the same reason; that can only drop a
frame that starts in the cycle after an acknowledgement freed a slot,
which the netlist would store, and never lands one on a slot still
held.
\item \textbf{The scanout fetches its own lines.} \code{scan}, a
\code{ScanFetch<32, 16, 640>}, takes each line request that arrives on
\code{scan\_req} from the pixel clock and starts \code{vfetch}, a
\code{LineFetch} behind \code{vhost}, whose words leave on
\code{scan\_words} for the line pair. The crossings are the board
top's, as the video slot's are. \code{vnobeats} holds the write side
the fetch never uses.
\item \textbf{The SD card host moves blocks through memory.}
\code{sdstore} and \code{sdfetch}, each behind a host of its own with
one bit of identifier, share the arbiter's sixth port through
\code{sdarb}, an \code{Arbiter2<32, 32, 4, 1, 2, 0>}. One command is a
read or a write and never both, and the host starts no command until
the last one's engine has finished, so the two never offer at once
and the nest costs neither a turn; the host checks that their busy
lines are never high together. \code{sdnoreads} and \code{sdnobeats}
hold the directions they never use.
\item \textbf{Razboj draws into the frame the scanout shows.}
\code{raster}, a \code{Raster<32, 2, 10, 480, 0x4200\_0000,
0x4280\_0000, 0x3900>} behind \code{rhost}, takes the arbiter's seventh
port. It reads its display list at \code{0x4280\_0000} and writes
rows of 1024 pixels from \code{0x4200\_0000}, the scanout's stride. Its
count is \code{doorbell}'s register on the tenth slot, which drives the
rasteriser's \code{ring} while the count is not zero, so an idle
rasteriser reads nothing; the join runs the rasteriser before the
doorbell, which reads its \code{idle} line.
\item The join runs the SPI master and the window before
\code{cfgflash}, since the pins read the lines both drive in the same
step. The flash's data out, \code{fl\_miso}, goes to both directly. The
flash's select and data in leave as \code{fl\_cs\_n} and
\code{fl\_mosi}; \code{fl\_cclk} is the startup block's model of the
clock pin, which a board leaves open, and \code{fl\_refused} says the
pins refused write enable. The master's interrupt is not wired.
\item \code{board\_netlist} puts a program in the core's
\code{imem} and in the data memory's \code{lane0} to \code{lane3} of
the netlist, since nothing on the machine loads them at run time.
\end{itemize}

The top, \code{vreteno\_board.v}, does the following around the
lowered \code{board}:

\begin{itemize}
\item It hands the board's 200~MHz differential clock to the memory
controller on \code{sys\_clk}, which makes the design's 100~MHz,
\code{ui\_clk}, from it; the board has no clock generator of its own.
\item The controller's reset, \code{sys\_rst}, is a pulse at power-on
counted on the board's clock and nothing else. The core's reset
follows the button, \code{ui\_rst} being low and \code{calib}, through
two flip-flops on the design's clock.
\item It ties \code{irq} low.
\item It puts the flash's select and data on the configuration pins
FCS\_B, D00 and D01, and holds D02 and D03, the chip's write protect
and hold, high. The clock, CCLK, is reached through \code{STARTUPE2}
inside the netlist.
\item Its four LEDs, lit when driven low, show the modulator's first
channel or the core halted, the serial line driven at least once, the
memory calibrated, and a heartbeat from bit 25 of a counter, which
says the controller's clock is good.
\end{itemize}

\subsection*{Verification}

\begin{itemize}
\item \code{//cpu/vreteno:board\_test} runs \code{Board<4>} in
Rust. \code{the\_greeting\_runs\_on\_the\_board} must print
\code{hello from rust} and halt within 8000 cycles.
\code{the\_memory\_test\_runs\_on\_the\_board} must print
\code{ddr3 ok} and halt within 40000 cycles.
\code{input\_comes\_by\_interrupt\_one\_byte\_each} runs a program
that takes four typed bytes through the interrupt controller, one
interrupt each, and must print \code{got ping in 4 interrupts} and
halt within 20000 cycles.
\code{a\_frame\_received\_lands\_in\_memory\_and\_the\_core\_reads\_it\_back}
puts a frame of 21 bytes on \code{net\_rx}, and the core must print its
length and every byte from the slot it landed in.
\code{two\_frames\_written\_to\_memory\_leave\_the\_port\_as\_written\_in\_order}
has the core write two frames of 23 and 18 bytes and ask for both back
to back, and the bytes that leave on \code{net\_tx} must be the two
frames, whole and in order.
\code{the\_entropy\_source\_answers\_on\_the\_board} must print
\code{trng ok}, place every pair of raw windows and find the run it
joins whole among the model's samples, within 120000 cycles.
\code{the\_configuration\_flash\_answers\_on\_the\_board} gives the
board a flash model holding a bitstream's first bytes; the core must
print the chip's identity, a clear write enable latch and where the
sync word is, and the chip must have taken no write enable.
The same target also runs the boot memory, the remote peripheral, the
loader and the debug module through the board.
\code{the\_netlist\_holds\_the\_controller} checks that the Verilog
has the top, the core and the DDR3's pins part as modules,
instantiates \code{ddr3\_axi32} and \code{ring\_osc} without writing
either, and has
the \code{dq} pads.
\item \code{//cpu/vreteno/board/sim:board\_test} simulates the top, the
netlist of \code{//cpu/vreteno:board\_sim\_v}, the controller and two
Micron models under Vivado's simulator. It is manual.
\code{//docs:vreteno} reports that in that run the core says
\code{ddr3 ok} on the serial port and halts.
\item \code{//cpu/vreteno:vreteno\_board\_synth} and
\code{//cpu/vreteno:vreteno\_board\_pnr} synthesise, place and route
the design with the board's pins. Both are manual.
\code{//docs:vreteno} reports that every timing constraint is met, with
0.307~ns worst setup slack and 0.033~ns worst hold slack, and 12\,296
LUTs, 11\,209 flip-flops, five block RAM tiles and four DSP blocks, for
the board at eed5b7f, with the configuration flash and the entropy
source's capture buffer; it has not been measured since.
\item \code{//cpu/vreteno:vreteno\_board\_prog} and
\code{//cpu/vreteno:vreteno\_board\_flash} program the FPGA and the
QSPI flash. Both are manual.
\code{//cpu/vreteno:vreteno\_board\_fits\_test}, also manual, checks
that the routed bitstream ends before \code{0x00A0\_0000}, where
programs begin in the flash.
\end{itemize}

No waveform target simulates \code{Board}'s netlist against a trace.

\subsection*{Used by}

\code{//cpu/vreteno:board\_netlist}, whose genrules
\code{//cpu/vreteno:board\_v} and \code{//cpu/vreteno:board\_sim\_v}
write the netlist, and \code{vreteno\_board.v}, which instantiates it
as \code{board}. The test \code{cpu/vreteno/tests/board.rs}.

\subsection*{Limits}

\code{//docs:vreteno} reports the design with the DDR3 memory and MIG
run on the board on September 28, 2026, written into the part over
JTAG, and saying \code{ddr3 ok}; the run from the flash after a power
cycle is owed (issue 188). It reports the configuration flash read on
the board as well: the identity, write enable refused and the
bitstream through the window. The program is fixed in the netlist. The board top ties \code{irq} low, so on the
board the interrupt controller's second source never asks.

The arbiter's ports are the core on 0, \code{jarb} on 1, the Ethernet
port's fetch engine on 2 and its store engine on 3, the scanout's fetch
on 4, \code{sdarb} on 5 and Razboj's rasteriser on 6. It takes turns,
so with every host offering each wins one turn in seven. Three bits of
port number leave room for an eighth host, which raises the count and
widens nothing past the arbiter.
